% =============================================================================
% topic_13.tex — Content only. No \documentclass.
% Compile as part of series via main.tex (one level up)
% =============================================================================

\setcounter{section}{0}

% ── Title page ────────────────────────────────────────────────────────────────
\maketitlepage
  {Topic 13}
  {Gravitational Fields}
  {Newton's Law \(\cdot\) Field Strength \(\cdot\) Potential \(\cdot\) Orbits \(\cdot\) Escape Velocity}
  {%
    \begin{itemize}[leftmargin=1cm, itemsep=2pt, label=\textbullet, topsep=4pt]
      \item Newton's Law of Gravitation
      \item Gravitational Field Strength
      \item Gravitational Potential
      \item Circular Orbits and Satellites
      \item Escape Velocity
      \item Formula Summary, Worked Examples, Exam Practice \& Checklist
    \end{itemize}
  }

% ── Table of contents ─────────────────────────────────────────────────────────
\tableofcontents
\newpage

% =============================================================================
\section{Core Concepts and Definitions}
% =============================================================================

\subsection{What is a Gravitational Field?}

A gravitational field is a region of space in which any object with mass experiences a
force. It is a \textbf{vector field} — at every point it has both magnitude and
direction (always directed towards the centre of mass of the body creating the field).

\subsection{Radial vs Uniform Fields}

\begin{itemize}
  \item \textbf{Radial field} (e.g.\ around a planet): field lines point towards the
        centre; field strength varies with distance. Represented by lines that converge
        at the centre.
  \item \textbf{Uniform field} (e.g.\ near Earth's surface over small distances):
        field lines are parallel and equally spaced; $g$ is approximately constant at
        $9.81\,\mathrm{N\,kg}^{-1}$.
\end{itemize}

% =============================================================================
\section{Newton's Law of Gravitation}
% =============================================================================

\begin{formulabox}{Newton's Law of Gravitation}
\begin{equation}
  F = \frac{Gm_1 m_2}{r^2}
\end{equation}
\begin{tabular}{ll}
$F$       & = gravitational force between the two masses (N) \\
$G$       & = gravitational constant $= 6.67\times10^{-11}\,\mathrm{N\,m^2\,kg^{-2}}$ \\
$m_1,m_2$ & = the two point masses or uniform spheres (kg) \\
$r$       & = separation between centres of mass (m) \\
\end{tabular}
\end{formulabox}

\subsection{Important Features of Newton's Law}
\begin{itemize}
  \item The force is always \textbf{attractive} — gravity never repels.
  \item It is an \textbf{inverse square law}: doubling $r$ reduces $F$ by a factor
        of 4.
  \item It applies strictly to point masses, but also works for uniform spheres if
        $r$ is measured from the centre.
  \item Newton's third law applies: both masses experience equal and opposite forces.
\end{itemize}

\begin{tipbox}
If asked to show that gravity obeys an inverse-square law, demonstrate that
$F \propto 1/r^2$. This might involve calculating $Fr^2$ for different values and
showing it is constant, or comparing the ratio $F_1/F_2 = r_2^2/r_1^2$.
\end{tipbox}

% =============================================================================
\section{Gravitational Field Strength}
% =============================================================================

\begin{defnbox}{Gravitational Field Strength $g$}
The gravitational field strength at a point is the gravitational force per unit mass
acting on a small test mass placed at that point.
\[
  g = \frac{F}{m} \qquad \text{units: N\,kg}^{-1}
\]
Note: $g$ is a vector quantity (directed towards the centre of mass of the source).
\end{defnbox}

\begin{formulabox}{Field Strength Equations}
\begin{equation}
  g = \frac{F}{m} \qquad \text{and} \qquad g = \frac{GM}{r^2}
\end{equation}
\begin{tabular}{ll}
$g$ & = gravitational field strength (N\,kg$^{-1}$ or m\,s$^{-2}$) \\
$F$ & = gravitational force on a test mass (N) \\
$m$ & = test mass (kg) \\
$M$ & = mass of body creating the field (kg) \\
$r$ & = distance from centre of $M$ (m) \\
\end{tabular}
\end{formulabox}

\begin{warnbox}
Do not confuse $g$ (field strength, N\,kg$^{-1}$) with $G$ (the universal
gravitational constant, N\,m$^2$\,kg$^{-2}$). They are completely different
quantities!
\end{warnbox}

% =============================================================================
\section{Gravitational Potential}
% =============================================================================

\begin{defnbox}{Gravitational Potential $\phi$}
The gravitational potential at a point is the work done per unit mass to move a small
test mass from infinity to that point.
\[
  \phi = \frac{W}{m} \qquad \text{units: J\,kg}^{-1}
\]
Gravitational potential is \textbf{always negative} (zero at infinity; work is done
by gravity as mass moves inwards, so the system loses potential energy).
\end{defnbox}

\begin{formulabox}{Gravitational Potential and Potential Energy}
\begin{equation}
  \phi = -\frac{GM}{r} \qquad \qquad E_p = m\phi = -\frac{GMm}{r}
\end{equation}
\begin{tabular}{ll}
$\phi$ & = gravitational potential at distance $r$ (J\,kg$^{-1}$) \\
$E_p$  & = gravitational potential energy (J) \\
$G$    & = $6.67\times10^{-11}\,\mathrm{N\,m^2\,kg^{-2}}$ \\
$M$    & = mass of body creating the field (kg) \\
$r$    & = distance from centre of $M$ (m) \\
\end{tabular}
\end{formulabox}

\subsection{Relationship Between $g$ and $\phi$}

The gravitational field strength $g$ is related to the potential $\phi$ by:

\begin{formulabox}{Field Strength as Potential Gradient}
\begin{equation}
  g = -\frac{\Delta\phi}{\Delta r}
\end{equation}
$g$ is the \textbf{negative gradient} of the potential–distance graph. A steeper
gradient means a stronger field.
\end{formulabox}

\begin{tipbox}
On a $\phi$ vs $r$ graph, the gradient at any point gives $-g$ at that point.
The graph curves towards zero as $r \to \infty$, and the steepest part is near
the surface where the field is strongest.
\end{tipbox}

\subsection{Equipotential Surfaces}

An \textbf{equipotential surface} is a surface on which the gravitational potential
is the same everywhere.
\begin{itemize}
  \item No work is done moving a mass along an equipotential surface.
  \item Equipotentials are always \textbf{perpendicular} to field lines.
  \item For a point mass (or uniform sphere), equipotentials are concentric spheres.
\end{itemize}

% =============================================================================
\section{Circular Orbits and Satellites}
% =============================================================================

For an object of mass $m$ in a circular orbit of radius $r$ around a mass $M$, the
gravitational force provides the centripetal force:

\begin{formulabox}{Orbital Mechanics}
\begin{equation}
  \frac{GMm}{r^2} = \frac{mv^2}{r} = mr\omega^2 = mr\left(\frac{2\pi}{T}\right)^2
\end{equation}
\begin{align}
  v &= \sqrt{\frac{GM}{r}} \qquad \text{(orbital speed)} \\[6pt]
  T^2 &= \frac{4\pi^2}{GM}\,r^3 \qquad \text{(Kepler's Third Law)}
\end{align}
\end{formulabox}

\begin{formulabox}{Kepler's Third Law}
\begin{equation}
  T^2 \propto r^3
\end{equation}
For all objects orbiting the same central mass $M$, the square of the orbital period
is proportional to the cube of the orbital radius.
\[
  \frac{T_1^2}{T_2^2} = \frac{r_1^3}{r_2^3}
\]
\end{formulabox}

\subsection{Energy in Circular Orbits}

\begin{center}
\renewcommand{\arraystretch}{1.5}
\begin{tabular}{lll}
\toprule
\textbf{Energy} & \textbf{Expression} & \textbf{Sign} \\
\midrule
Kinetic energy       & $E_k = \dfrac{GMm}{2r}$  & Positive \\[6pt]
Gravitational PE     & $E_p = -\dfrac{GMm}{r}$  & Negative \\[6pt]
Total energy         & $E_{\text{total}} = -\dfrac{GMm}{2r}$ & Negative \\
\bottomrule
\end{tabular}
\end{center}

\medskip
Note: $E_{\text{total}} = \tfrac{1}{2}E_p$. Total energy is always negative for a
bound orbit.

\begin{warnbox}
When a satellite moves to a \textbf{lower} orbit (smaller $r$): its speed
\emph{increases}, its kinetic energy \emph{increases}, but its total energy
\emph{decreases} (becomes more negative). This is counterintuitive — the satellite
speeds up yet loses total energy because the drop in potential energy outweighs the
gain in kinetic energy.
\end{warnbox}

\subsection{Geostationary Orbits}

\begin{defnbox}{Geostationary Satellite}
A geostationary satellite has the following properties:
\begin{itemize}[nosep]
  \item Orbital period $T = 24\,\text{hours}$ (same as Earth's rotation)
  \item Orbits directly above the equator
  \item Orbits in the same direction as Earth's rotation (west to east)
  \item Appears stationary relative to a point on Earth's surface
  \item Orbital radius $\approx 42\,000\,\text{km}$ from Earth's centre
  \item Used for: telecommunications, weather satellites, GPS
\end{itemize}
\end{defnbox}

% =============================================================================
\section{Escape Velocity}
% =============================================================================

\begin{defnbox}{Escape Velocity}
The \textbf{escape velocity} is the minimum speed needed for an object to escape from
a gravitational field to infinity, without further propulsion.
\end{defnbox}

\begin{formulabox}{Escape Velocity}
\begin{equation}
  v_{\text{esc}} = \sqrt{\frac{2GM}{R}}
\end{equation}
Derived by setting total energy to zero at infinity:
\[
  \frac{1}{2}mv^2 - \frac{GMm}{R} = 0
  \quad\Longrightarrow\quad
  v_{\text{esc}} = \sqrt{\frac{2GM}{R}}
\]
Note: escape velocity is \textbf{independent} of the mass of the escaping object.
\end{formulabox}

\begin{tipbox}
The derivation of escape velocity is a common exam request. Set $E_k + E_p = 0$
at the point of just escaping (total energy $= 0$ at infinity), then solve for $v$.
\end{tipbox}

% =============================================================================
\section{Formula Summary Sheet}
% =============================================================================

\begin{summarybox}
\renewcommand{\arraystretch}{1.7}
\begin{tabular}{p{5.5cm} p{5.5cm} p{2cm}}
\toprule
\textbf{Formula} & \textbf{Meaning} & \textbf{Section} \\
\midrule
$F = \dfrac{Gm_1m_2}{r^2}$
  & Newton's Law of Gravitation & 2 \\[6pt]
$g = \dfrac{F}{m}$
  & Field strength (general) & 3 \\[6pt]
$g = \dfrac{GM}{r^2}$
  & Field strength (radial) & 3 \\[6pt]
$\phi = -\dfrac{GM}{r}$
  & Gravitational potential & 4 \\[6pt]
$E_p = -\dfrac{GMm}{r}$
  & Gravitational potential energy & 4 \\[6pt]
$g = -\dfrac{\Delta\phi}{\Delta r}$
  & Field strength as potential gradient & 4 \\[6pt]
$v = \sqrt{GM/r}$
  & Orbital speed & 5 \\[6pt]
$T^2 = \dfrac{4\pi^2 r^3}{GM}$
  & Kepler's Third Law & 5 \\[6pt]
$v_{\text{esc}} = \sqrt{2GM/R}$
  & Escape velocity & 6 \\
\midrule
\multicolumn{3}{l}{\small $G = 6.67\times10^{-11}\,\mathrm{N\,m^2\,kg^{-2}}$,\quad
  $M_{\oplus} = 5.97\times10^{24}\,\mathrm{kg}$,\quad
  $R_{\oplus} = 6.37\times10^{6}\,\mathrm{m}$} \\
\bottomrule
\end{tabular}
\end{summarybox}

% =============================================================================
\section{Exam Technique and Problem-Solving Strategy}
% =============================================================================

\subsection{5-Step Strategy for Gravitational Field Problems}
\begin{enumerate}
  \item \textbf{Identify the type of problem}: is it force, field strength, potential,
        orbit, or energy?
  \item \textbf{Draw a diagram}: mark all masses, distances from centre (not surface!),
        and directions.
  \item \textbf{Select the correct formula}: radial field ($GM/r^2$) or near-surface
        constant $g$?
  \item \textbf{Check units and powers of 10}: distances in metres, masses in kg,
        $r$ always from the centre.
  \item \textbf{Check sign conventions}: potential and potential energy are always
        negative; force and field strength magnitudes are positive.
\end{enumerate}

\subsection{Most Common Mistakes}
\begin{itemize}
  \item Using the radius of a planet as the surface height rather than measuring
        from the centre.
  \item Forgetting the negative sign on gravitational potential
        ($\phi = -GM/r$).
  \item Confusing gravitational potential (J\,kg$^{-1}$) with gravitational
        potential energy (J).
  \item In Kepler's Third Law, forgetting to convert $T$ to seconds or $r$ to
        metres.
  \item Applying centripetal force equations to objects on a surface (only valid
        for orbiting objects).
\end{itemize}

% =============================================================================
\section{Worked Examples}
% =============================================================================

\noindent\textbf{Question 1}\quad Calculate the gravitational field strength at a
height of 400\,km above Earth's surface.
$(M_E = 5.97\times10^{24}\,\mathrm{kg},\ R_E = 6.37\times10^6\,\mathrm{m})$

\begin{solutionbox}
$r = R_E + h = 6.37\times10^6 + 4.00\times10^5 = 6.77\times10^6\,\mathrm{m}$
\[
  g = \frac{GM}{r^2}
    = \frac{6.67\times10^{-11} \times 5.97\times10^{24}}{(6.77\times10^6)^2}
    = \frac{3.98\times10^{14}}{4.58\times10^{13}}
    = 8.69\,\mathrm{N\,kg^{-1}}
\]
Note: this is the field strength at the ISS orbit — astronauts are not weightless
because gravity is zero!
\end{solutionbox}

\bigskip
\noindent\textbf{Question 2}\quad A satellite orbits Earth at a radius of
$4.2\times10^7\,\mathrm{m}$. Calculate its orbital period.

\begin{solutionbox}
\[
  T^2 = \frac{4\pi^2 r^3}{GM}
      = \frac{4\pi^2 \times (4.2\times10^7)^3}{6.67\times10^{-11}
        \times 5.97\times10^{24}}
      = \frac{4\pi^2 \times 7.41\times10^{22}}{3.98\times10^{14}}
      = 7.36\times10^9\,\mathrm{s^2}
\]
\[
  T = \sqrt{7.36\times10^9} = 8.58\times10^4\,\mathrm{s} \approx 23.8\,\text{hours}
\]
This is a geostationary orbit ($T \approx 24\,\text{h}$).
\end{solutionbox}

\bigskip
\noindent\textbf{Question 3}\quad A spacecraft of mass 2500\,kg is moved from
Earth's surface to an altitude of 800\,km. Calculate the increase in gravitational
potential energy.

\begin{solutionbox}
$r_1 = 6.37\times10^6\,\mathrm{m}$,\quad
$r_2 = 6.37\times10^6 + 8.00\times10^5 = 7.17\times10^6\,\mathrm{m}$
\[
  \Delta E_p = -\frac{GMm}{r_2} - \left(-\frac{GMm}{r_1}\right)
             = GMm\left(\frac{1}{r_1} - \frac{1}{r_2}\right)
\]
\[
  \Delta E_p = 6.67\times10^{-11} \times 5.97\times10^{24} \times 2500
               \times \left(\frac{1}{6.37\times10^6} - \frac{1}{7.17\times10^6}\right)
             = 1.74\times10^{10}\,\mathrm{J}
\]
\end{solutionbox}

\newpage
% =============================================================================
\section{Practice Exam Questions}
% =============================================================================

\subsection*{Section A — Short Answer Questions}

\noindent\begin{minipage}{\textwidth}
\textbf{Q1.} State Newton's Law of Gravitation.
\markalloc{2}

\answerbox{2.5cm}
\end{minipage}

\bigskip
\noindent\begin{minipage}{\textwidth}
\textbf{Q2.} Define gravitational field strength and state its units.
\markalloc{2}

\answerbox{2.5cm}
\end{minipage}

\bigskip
\noindent\begin{minipage}{\textwidth}
\textbf{Q3.} Explain why gravitational potential is always a negative quantity.
\markalloc{2}

\answerbox{3cm}
\end{minipage}

\bigskip
\noindent\begin{minipage}{\textwidth}
\textbf{Q4.} The gravitational field strength at Earth's surface is
$9.81\,\mathrm{N\,kg}^{-1}$. Calculate the field strength at a distance of $3R_E$
from Earth's centre, where $R_E$ is Earth's radius.
\markalloc{2}

\answerbox{3cm}
\end{minipage}

\bigskip
\noindent\begin{minipage}{\textwidth}
\textbf{Q5.} Two planets, X and Y, orbit the same star. Planet X has an orbital
radius twice that of planet Y. Determine the ratio $T_X/T_Y$.
\markalloc{3}

\answerbox{4cm}
\end{minipage}

\subsection*{Section B — Longer Structured Questions}

\noindent\begin{minipage}{\textwidth}
\textbf{Q6.} A satellite of mass $m$ orbits a planet of mass $M$ in a circular orbit
of radius $r$.

\begin{enumerate}[label=(\alph*)]
  \item Show that the orbital speed of the satellite is given by
        $v = \sqrt{GM/r}$.
        \markalloc{2}

        \answerbox{3.5cm}

  \item Show that the period of the orbit satisfies
        $T^2 = \dfrac{4\pi^2 r^3}{GM}$.
        \markalloc{2}

        \answerbox{3.5cm}

  \item The satellite is moved to a lower orbit. Explain what happens to its
        speed, kinetic energy, and total energy.
        \markalloc{3}

        \answerbox{4cm}
\end{enumerate}
\end{minipage}

\bigskip
\noindent\begin{minipage}{\textwidth}
\textbf{Q7.} The Moon orbits the Earth with a period of 27.3 days at a mean orbital
radius of $3.84\times10^8\,\mathrm{m}$.

\begin{enumerate}[label=(\alph*)]
  \item Use this data to calculate the mass of the Earth.
        \markalloc{3}

        \answerbox{4cm}

  \item Calculate the gravitational potential at the Moon's orbital radius.
        \markalloc{2}

        \answerbox{3cm}

  \item A 75\,kg astronaut travels from Earth's surface to the Moon's orbital
        radius. Calculate the change in gravitational potential energy.
        \markalloc{3}

        \answerbox{4cm}
\end{enumerate}
\end{minipage}

\newpage
% =============================================================================
\section{Mark Scheme and Answers}
% =============================================================================

\begin{markbox}
\textbf{Q1.} Any two masses exert an attractive force on each other [1]; the force is
proportional to the product of their masses and inversely proportional to the square
of their separation [1].

\medskip
\textbf{Q2.} Gravitational field strength is the gravitational force per unit mass
acting on a (small test) mass placed at that point [1]; units: N\,kg$^{-1}$
(or m\,s$^{-2}$) [1].

\medskip
\textbf{Q3.} Gravitational potential is defined as zero at infinity [1]; work is done
by gravity as mass moves from infinity inward, so the potential decreases below zero —
it is always negative [1].

\medskip
\textbf{Q4.} $g \propto 1/r^2$; at $3R_E$:
$g = 9.81/3^2 = 9.81/9 = 1.09\,\mathrm{N\,kg^{-1}}$ [2].

\medskip
\textbf{Q5.} By Kepler's Third Law: $T^2 \propto r^3$, so
$\left(\dfrac{T_X}{T_Y}\right)^2 = \left(\dfrac{2r_Y}{r_Y}\right)^3 = 8$ [2];\quad
$\dfrac{T_X}{T_Y} = \sqrt{8} = 2\sqrt{2} \approx 2.83$ [1].

\medskip
\textbf{Q6(a).} Gravitational force provides centripetal force:
$\dfrac{GMm}{r^2} = \dfrac{mv^2}{r}$ [1]; simplify to give $v = \sqrt{GM/r}$ [1].

\textbf{Q6(b).} Substitute $v = 2\pi r/T$ into $v^2 = GM/r$:
$\dfrac{4\pi^2 r^2}{T^2} = \dfrac{GM}{r}$ [1];
rearrange to $T^2 = \dfrac{4\pi^2 r^3}{GM}$ [1].

\textbf{Q6(c).} Speed: increases (since $v \propto 1/\sqrt{r}$) [1]; kinetic energy:
increases [1]; total energy: decreases (becomes more negative, since
$E = -GMm/2r$) [1].

\medskip
\textbf{Q7(a).}
$T = 27.3 \times 24 \times 3600 = 2.36\times10^6\,\mathrm{s}$;\quad
$M = \dfrac{4\pi^2 r^3}{GT^2}
   = \dfrac{4\pi^2 (3.84\times10^8)^3}{6.67\times10^{-11}(2.36\times10^6)^2}$
[1] $= 6.02\times10^{24}\,\mathrm{kg}$ [2].

\textbf{Q7(b).}
$\phi = -\dfrac{GM}{r}
      = -\dfrac{6.67\times10^{-11} \times 6.02\times10^{24}}{3.84\times10^8}
      = -1.05\times10^6\,\mathrm{J\,kg^{-1}}$ [2].

\textbf{Q7(c).}
$\Delta E_p = m\Delta\phi = m(\phi_{\text{Moon orbit}} - \phi_{\text{surface}})$ [1];\quad
$\phi_{\text{surface}} = -GM/R_E = -6.25\times10^7\,\mathrm{J\,kg^{-1}}$;\quad
$\Delta\phi = -1.05\times10^6 - (-6.25\times10^7) = 6.14\times10^7\,\mathrm{J\,kg^{-1}}$;\quad
$\Delta E_p = 75 \times 6.14\times10^7 = 4.6\times10^9\,\mathrm{J}$ [2].
\end{markbox}

\newpage
% =============================================================================
\section{Revision Checklist}
% =============================================================================

\begin{checklistbox}
\small
\renewcommand{\arraystretch}{1.5}
\begin{tabular}{p{11cm} >{\centering\arraybackslash}p{2.5cm}}
\toprule
\textbf{Learning Objective} & \textbf{Confidence (1--3)} \\
\midrule
$\square$\; State and apply Newton's Law of Gravitation & \\
$\square$\; Define gravitational field strength; use $g = F/m$ and $g = GM/r^2$ & \\
$\square$\; Sketch field line and equipotential diagrams & \\
$\square$\; Define gravitational potential and explain why it is negative & \\
$\square$\; Use $\phi = -GM/r$ and $E_p = -GMm/r$ & \\
$\square$\; Apply the relationship $g = -\Delta\phi/\Delta r$ & \\
$\square$\; Derive the expression for orbital speed & \\
$\square$\; State and apply Kepler's Third Law ($T^2 \propto r^3$) & \\
$\square$\; Analyse energy changes in satellite orbits & \\
$\square$\; Describe the properties of a geostationary satellite & \\
$\square$\; Derive and apply the formula for escape velocity & \\
$\square$\; Interpret graphs of $g$ vs $r$ and $\phi$ vs $r$ & \\
\bottomrule
\end{tabular}

\medskip
\begin{center}
\textit{Key:\quad 1 = Need more work \quad 2 = Getting there \quad 3 = Confident}
\end{center}
\end{checklistbox}

\vspace{1cm}
\begin{center}
{\large\bfseries\color{cieprimary} Good luck with your revision!}\\[0.4cm]
\textit{Remember: understanding \emph{why} formulae work is more powerful than
memorising them.}\\[0.2cm]
\textit{Practise deriving orbital speed, Kepler's Third Law, and escape velocity
from first principles.}
\end{center}
